% Short sibling of ai_core_nontech_prompts.tex (comment-sync pattern): same frames, only the ones the short seminar needs are live.
% Mirror any frame added, removed or edited in ai_core_nontech_prompts.tex here. Trim by commenting frames, not deleting.

%%%%%%%%%%%%%%%%%%%%%%%%%%%%%%%%%%%%%%%%%%%%%%%%%%%%%%%%%%%
\begin{frame}[fragile]\frametitle{What Is a Prompt?}
\begin{itemize}
\item A \textbf{prompt} is the text you type to the AI: a question, an instruction, or both.
\item The model reads your prompt and continues it, one word at a time, just like the SMS suggestion.
\item So your words steer the answer: more specific text gives a more specific result.
\item A prompt can hold a role for the AI, background, an example, and the text to work on.
\item \textbf{Prompt engineering} is writing prompts that give usable answers.
\end{itemize}
\end{frame}

%%%%%%%%%%%%%%%%%%%%%%%%%%%%%%%%%%%%%%%%%%%%%%%%%%%
\begin{frame}[fragile]\frametitle{Same Task, Different Prompts}
\begin{itemize}
\item Same AI, same topic: a vague prompt gives a generic answer, a specific one gives a usable answer.
\item A chat remembers what came before. Reply to refine the answer instead of starting again.
\end{itemize}
\begin{lstlisting}[language={}]
Vague:      Write a notice about water tank cleaning.

Specific:   Write a notice for a housing society about
            water tank cleaning on Sunday, 6 to 10 am,
            in 80 words, polite tone.

Follow-up:  Make it shorter, and add a contact number
            placeholder.
\end{lstlisting}
\end{frame}

%%%%%%%%%%%%%%%%%%%%%%%%%%%%%%%%%%%%%%%%%%%%%%%%%%%%%%%%%%%
\begin{frame}[fragile]\frametitle{A Framework for Every Prompt: CRAFT}

\begin{center}
\Large \textbf{C R A F T}
\end{center}

\medskip
\begin{itemize}
    \item \textbf{C}ontext, who you are, what the situation is
    \item \textbf{R}ole, who the AI should act as
    \item \textbf{A}ction, exactly what you want it to do
    \item \textbf{F}ewshots, one or two examples of what ``good'' looks like
    \item \textbf{T}weaks, tone, length, format, and what to avoid
\end{itemize}

\medskip
\begin{center}
\textbf{Run through the five letters in your head before you type. That is the whole framework.}
\end{center}

\end{frame}

%%%%%%%%%%%%%%%%%%%%%%%%%%%%%%%%%%%%%%%%%%%%%%%%%%%%%%%%%%%
% \begin{frame}[fragile]\frametitle{CRAFT: Context and Role}
%
% \begin{itemize}
%     \item \textbf{Context}, without it, the AI is guessing: what the subject is, who the reader is, what has already happened.
%     \item \textbf{Role}, telling the AI who to act as changes its whole register: a ``friendly neighbour'' writes differently from a ``formal clerk.''
% \end{itemize}
%
% \begin{lstlisting}[language=HTML]
% Weak:  Write a notice about water tank cleaning.
%
% CRAFT: You are the secretary of a housing society (Role).
%        The tank is cleaned twice a year, and residents
%        have complained about short notice before (Context).
% \end{lstlisting}
%
% \end{frame}

%%%%%%%%%%%%%%%%%%%%%%%%%%%%%%%%%%%%%%%%%%%%%%%%%%%%%%%%%%%
% \begin{frame}[fragile]\frametitle{CRAFT: Action and Fewshots}
%
% \begin{itemize}
%     \item \textbf{Action}, one clear instruction. ``Help me with X'' is not an action, ``Write a 3-line notice about X'' is.
%     \item \textbf{Fewshots}, show, do not just tell. One example of the style or format you want is worth three sentences of description.
% \end{itemize}
%
% \begin{lstlisting}[language=HTML]
% Action:   Write a 3-line notice about the tank
%           cleaning on Sunday.
%
% Fewshot:  Match this style:
%           "Water supply will be off on Monday,
%            9 to 12. Please store water in advance."
% \end{lstlisting}
%
% \end{frame}

%%%%%%%%%%%%%%%%%%%%%%%%%%%%%%%%%%%%%%%%%%%%%%%%%%%%%%%%%%%
% \begin{frame}[fragile]\frametitle{CRAFT: Tweaks}
%
% \begin{itemize}
%     \item \textbf{Tweaks} are the small dials that turn a generic answer into a usable one:
%     \begin{itemize}
%         \item \textbf{Tone}, formal / friendly / encouraging
%         \item \textbf{Length}, word count or ``under 5 sentences''
%         \item \textbf{Format}, bullet list / table / email / paragraph
%         \item \textbf{Constraints}, what to explicitly avoid
%     \end{itemize}
%     \item Add tweaks last, once Context, Role, Action, and Fewshots already give the AI enough to work with.
% \end{itemize}
%
% \begin{lstlisting}[language=HTML]
% Tweaks: Keep it under 60 words, friendly tone,
%         no exclamation marks.
% \end{lstlisting}
%
% \end{frame}

%%%%%%%%%%%%%%%%%%%%%%%%%%%%%%%%%%%%%%%%%%%%%%%%%%%%%%%%%%%
\begin{frame}[fragile]\frametitle{CRAFT, All Together}

\begin{lstlisting}[language=HTML]
You are the secretary of a housing society (R).
The tank is cleaned twice a year, and residents
have complained about short notice before (C).

Write a 3-line notice about the tank cleaning
on Sunday, 6 to 10 am (A).

Match this style: "Water supply will be off on
Monday, 9 to 12. Please store water in advance." (F)

Keep it under 60 words, friendly tone,
no exclamation marks (T).
\end{lstlisting}

\end{frame}

%%%%%%%%%%%%%%%%%%%%%%%%%%%%%%%%%%%%%%%%%%%%%%%%%%%%%%%%%%%
\begin{frame}[fragile]\frametitle{CRAFT: The Takeaways}

\begin{itemize}
    \item Not every prompt needs all five letters spelled out explicitly, but running through them in your head catches what you would otherwise forget.
    \item Save a prompt that works. A small personal library of tested prompts saves more time than any single clever prompt.
    \item \textbf{Always remember:} treat the output as a first draft, review it, and never paste confidential data into a public AI tool.
\end{itemize}

\end{frame}
